% NAL 2338, batch 1 — Gallica views 1–20.
% Pass: Opus 5.5 (claude-opus-5-5), 2026-09-23 — first pass, unchecked against the views by a human.
%
% What the twenty views hold. Greyscale scans (seven in colour: v. 1–3, 8,
% 10, 14, 18; all read in grey), portrait, about 4316 × 5740, one page per view, the leaf lying on the
% opened volume with the edge of the next leaf or of the guards showing at one
% side. Nineteen views carry text. Five pieces, all Latin letters to or about
% Mersenne and Roberval on the cycloid, telescope glasses and the quadratures,
% in three hands:
%
%   v. 1       The library's title leaf: « Lat. nouv. Acq. 2338 », « Papiers
%              de Mersenne », « Volume de 50 Feuillets / 1er Octobre 1888. »,
%              « Libri 1846. ». Not transcribed.
%   v. 2–7     PIECE 1. Letter headed « M. Mers:no Clarissmo, et Celeberrimo
%              Viro S.P. », dated « Dat: Flor: Kal: Maij: anno 1644 », signed
%              « Addictissmus Seruus / Euanga Torricellius. » (v. 6), address
%              « Celeberrimo Viro P. Marino Mersenno / Lut. Paris. » on the
%              last verso (v. 7). Hand A. Galileo left almost nothing on the
%              force of percussion; Roberval's rectification of the parabola;
%              du Verdus; the cycloid (« Cycloides secundarias »,
%              « genitor ») with a small figure (v. 3): area of the cycloidal
%              space, tangents, solids of revolution about the tangent at the
%              vertex (subsesquitertium, credited to Antonio Nardi), about a
%              tangent parallel to the base, about the axis (11 to 18), about
%              the base (5 to 8, credited to Roberval « ante nos »); centre of
%              gravity of the cycloid (7 to 5) and of the semicycloid (19 and
%              8; 16 to 11); Doni; the Grand Duke's glass works; Mersenne's
%              Cogitata physico-mathematica awaited.
%   v. 8–11    PIECE 2. « Doctissimo et celeberrimo Viro P. M. Mersenno
%              Euangta: Torricellius S. », ending « Flor: die 7 Iulij 1646.
%              Euangelista Torricellius. » and a postscript (v. 11). Hand B,
%              a posed upright italic, a fair copy or a secretary's hand: the
%              signature is in the same hand and module as the body. Letters of
%              Mersenne and Roberval dated Kal. Ian. received at the end of
%              March through Ferdinando Bardi and Andrea Arrighetti; the
%              quality of Torricelli's lenses against Fontana's, Francini's
%              and Galileo's (the Gassendi tube); a ten-« passus » telescope
%              made from a galley mast; Baliani's examen of a quadrature by
%              « Batanus »; the quadrature from « Ars Magna lucis,
%              et Vmbræ » (« à Mathematico Romano Societatis
%              Iesu »), with marginal page references « Pagina 317 », « Pagina
%              327 versu 10 »; spirals for Roberval, hyperbolas for Fermat, a
%              « Carollariolum » announced but not on the leaves.
%   v. 12–15   PIECE 3. « Clmo: Viro P. Marino Mersenno / Euangta Torricellius
%              S. D. », ending « Vale. / Florentiæ Kalendis Februarij anno
%              1647. » without subscription (the name stands in the last
%              sentence). Hand A again. Roberval's delayed answer; the centre
%              of percussion (« Problema … non arridet animo meo »); the
%              concave glass; the drawn circle giving the aperture (v. 13);
%              Jupiter's belts; Torricelli's spirals and infinite hyperbolas;
%              Roberval's announced epistle; the reprint of his « libellos »;
%              Fermat's problem of the point minimising the sum of distances
%              to three given points, « construxi, demonstraui,
%              determinauique »; Cavalieri, Magiotti, Renieri; Niceron's book
%              bequeathed.
%   v. 16–19   PIECE 4. « Doctissime, et Celeberrime P. Mersenne », a
%              letter in the first person, unsigned and undated, ending
%              « Atque hic me tibi iterum commendo / Torricellium » (v. 19).
%              Hand C, a heavy sloping cursive. It quotes under its own
%              heading (« Incomparabili Geometræ D. Torricellio S.P.D. ») the
%              chapters of a letter from Mersenne received « sub die 28 Iunij
%              1644 » and the end of the same letter (« uersus finem iterum
%              hæc habet P. Vra »), to claim priority against Roberval for the
%              centre of gravity of the cycloid and the solid about the axis.
%              Boxed notes in a smaller module (« uerbum deest », « Parata bis
%              legitur », v. 16; « aduit tandem licet breui tempore », v. 18)
%              are a copyist's remarks, which makes the piece a copy; whose,
%              the leaves do not say.
%   v. 20      PIECE 5, beginning. « Clarmo Viro Roberuallo / Euangta
%              Torricellius S.P. », hand B. The initial and the first letters
%              of the first six lines have faded or were never finished. A
%              later hand wrote « 1 Sept. 1643 » at the top left; the letter
%              itself carries no date on this page. The centre of gravity of
%              the parabola; « Mensura Cycloidis (hoc enim nomine Clar:
%              Galileus appellauit 45 iam ab hinc annis figuram …) ».
%
% Continuity. View 1 is the title leaf, so nothing runs in from before. View
% 20 does NOT end its piece: its last line, « Illam deinde quinquies diuersis
% semper principijs », runs on into view 21.
%
% Who writes to whom. Pieces 1–3 and 5 are headed with Torricelli's name and
% addressed to Mersenne (1–3) and to Roberval (5); piece 1 alone carries an
% autograph-looking signature and an address panel. Piece 4 is addressed to
% Mersenne and closes on « Torricellium » as the object of « commendo »; its
% writer is not otherwise named. Nothing on these leaves is in a hand that the
% leaves themselves attribute to Mersenne. The dates on the leaves: Kal. Maij
% 1644 (v. 6), 7 Iulij 1646 (v. 11), Kal. Februarij 1647 (v. 15), « sub die 28
% Iunij 1644 » quoted (v. 16), « Kal. Ian. » (v. 8, of letters received), and
% the later « 1 Sept. 1643 » (v. 20).
%
% Three hands.
% (A) v. 2–6, 12–15: a fast Italian cursive, sloping, with large looped
%     capitals. The final m after a vowel is often reduced to a descending
%     loop (etiã-like); it is set as m throughout, not as a suspension. The
%     e-caudata (ę) is set æ. -que as q; is kept. The double s looks like n
%     (« fecisses »). Accents on adverbs and ablatives (adeò, penè, repetitâ,
%     lectionê) are kept.
% (B) v. 8–11, 20: a posed upright italic, almost without correction, with
%     long line-fillers at line ends (not transcribed); writes æ, v initial
%     and u medial (vitrum, vniuersis, reseruauisse); p with a crossed
%     descender for per, set ꝑ; q; kept.
% (C) v. 16–19: a heavy sloping cursive whose final t is an 8-shaped loop
%     (erat, habebat, et, putet) and whose 8 in numbers is half-formed and can
%     be read 0 (« 11 ad 18 », « 11 et 18 », « 5 ad 8 », all flagged).
% All three write u medially where a modern text has v (« Roberuallius »,
% « breuitate »); the leaves' u and v are kept.
%
% Notation. None beyond letters of figures and ratios of integers written in
% words (« ut 7 ad 5 »). Greek once, « ἀδύνατον » (v. 16), kept in Greek. A
% long s is transcribed as s. Figures: v. 3 (circle AB, radius CAE, the
% cycloid AD dotted) and v. 13 (a plain circle, the aperture), described in
% notes.
%
% Numbers on the leaves.
% - One thin ink series at the top right of every recto, across the three
%   hands: 1 (v. 2), 2 (v. 4), 3 (v. 6), 4 (v. 8), 5 (v. 10), 6 (v. 12), 7
%   (v. 14), 8 (v. 16), 9 (v. 18), 10 (v. 20). No verso carries a number; the
%   title leaf (v. 1) is not counted. It behaves as the library's foliation of
%   the 1888 certificate (« Volume de 50 Feuillets »), but it is ink, not
%   pencil, so, as for NAF 5176, NAF 9544 and Français 12357, no \folio{} is
%   written until a human has looked. Recorded in the notes of each recto.
% - « R.c. 8070 (86) » at the foot of v. 2, with the oval stamp « ACQ. 8070
%   (LIBRI) »: the acquisition mark. Round stamps « BIBLIOTHÈQUE NATIONALE
%   MANUSCRITS / R.F. » on v. 2, 8, 12, 16, 20.
% - « Pagina 317 », « Pagina 327 versu 10 » (v. 10, left margin, hand B): the
%   writer's references to the pages of the book he cites, not leaf numbers.
% - « 1 Sept. 1643 » (v. 20): a later dating in another hand, not a number
%   of the series.
% - No writer's pagination.
%
% The catalogue gives no dating: \dating{} is left empty, and this pass
% infers no date for the volume. The dates above are the leaves'.

% Coordinator's reconciliation (2026-09-23), after all six batches:
%   One ink series of leaf numbers runs top right of the rectos across every
%   hand and piece: 1–10 (v. 2–20), 11–20 (v. 22–40), 21–29 (v. 42–59),
%   30–38 with « 34 bis » for the figure slip (v. 61–79), 39–48 (v. 81–99),
%   49–50 (v. 101–103). It ends on « 50 » at the last written leaf, as the
%   certificate of v. 1 has it (« Volume de 50 Feuillets / 1er Octobre
%   1888 »): it is the library's count. Being ink, it gets no \folio{}, in any
%   batch. The bold series on the letter of v. 24–35 (1–11) is that letter's
%   own pagination (batch 2).
%   Seams: v. 60 ends « vniuersa- » and v. 61 opens « lius »; batch 3 had set
%   the catchword into v. 60, and now ends where the page does. v. 100 runs
%   into v. 101 (« altius ; » / « in aliis autem humilius »), the « De Vacuo
%   Narratio » (v. 97–103). The « non » sign is set « nl » in batch 1 (hand B,
%   an n with a barred ascender) and « n^o » in batch 2 (the fast cursive, an
%   n with a raised loop): two hands, two signs, left as each pass read them.

\documentclass[11pt,a4paper]{article}
% The preamble shared by every transcription of the mathematicians' manuscripts in Gallica.
%
% It rests on one idea: **uncertainty compounds, it does not resolve.** A
% transcription that smooths over a doubtful word has destroyed the only thing
% it was made for — a reader must be able to tell, at every point, what was on
% the page and what was supplied. The apparatus macros below are therefore the
% critical apparatus, not decoration.
%
% The same commands are recognised by `scripts/render.mjs`, which produces the
% site's reading view, and by `scripts/tei.mjs`. A macro added here must be
% added there too, or rendering fails loudly — which is the wanted behaviour.
%
% Comments here are English, like the rest of the repository. The documents
% this preamble sets are French, and that is deliberate: see the skills.

\ProvidesPackage{germain}[2026/09/15 Transcriptions of mathematicians' manuscripts in Gallica]

\usepackage{amsmath,amssymb,amsthm}
\usepackage{mathrsfs}
% Kept from the parent preamble so the renderer's diagram support stays
% available; these manuscripts draw no commutative diagrams, and nothing here uses it.
\usepackage{tikz-cd}
\usepackage[english,french]{babel}
\usepackage{fontspec}

% --- Greek ------------------------------------------------------------------
% The text face stays fontspec's default, Latin Modern, which has no Greek:
% XeTeX drops every glyph it lacks with a line in the log and nothing on the
% page, and the Greek words the manuscripts quote vanished from the PDFs. So
% the Greek and Greek Extended blocks get a XeTeX character class of their own,
% and entering or leaving a run of it switches the family to CMU Serif — the
% same Computer Modern design, with polytonic Greek — and back. Only the
% family changes: a Greek word in \textit or \textbf keeps its shape and series.
% The transitions to and from every other class carry the tokens that class
% has towards an ordinary letter, so French spacing before « ; » or inside
% « » still holds after a Greek word. No grouping, so no brace or \endgroup
% can come out unbalanced; maths is left alone, since XeTeX inserts nothing
% there. Kept identical in `exercices.sty`.
\makeatletter
\newfontfamily\germainGreekfont{cmun}[
  Extension=.otf, NFSSFamily=germaingreek,
  UprightFont=*rm, ItalicFont=*ti, SlantedFont=*sl,
  BoldFont=*bx, BoldItalicFont=*bi, BoldSlantedFont=*bl]
\newXeTeXintercharclass\germain@greek
\def\germain@greekrange#1#2{%
  \@tempcnta=#1\relax
  \loop
    \XeTeXcharclass\@tempcnta=\germain@greek
    \ifnum\@tempcnta<#2\advance\@tempcnta\@ne\repeat}
\germain@greekrange{"0370}{"03FF}
\germain@greekrange{"1F00}{"1FFF}
\def\germain@greekprev{\rmdefault}
\def\germain@greekin{%
  \edef\germain@greekprev{\f@family}\fontfamily{germaingreek}\selectfont}
\def\germain@greekout{\fontfamily{\germain@greekprev}\selectfont}
\def\germain@greekjoin{%
  \XeTeXinterchartokenstate=\@ne
  \@tempcnta=\z@
  \loop
    \ifnum\@tempcnta=\germain@greek\else
      \XeTeXinterchartoks\germain@greek\@tempcnta=\expandafter{%
        \expandafter\germain@greekout\the\XeTeXinterchartoks\z@\@tempcnta}%
      \XeTeXinterchartoks\@tempcnta\germain@greek=\expandafter{%
        \the\XeTeXinterchartoks\@tempcnta\z@\germain@greekin}%
    \fi
    \ifnum\@tempcnta<4095\advance\@tempcnta\@ne\repeat}
% babel-french sets its own classes, and saves and restores the interchar
% state, each time a language is selected: join again after it does.
\addto\extrasfrench{\germain@greekjoin}
\addto\extrasenglish{\germain@greekjoin}
\AtBeginDocument{\germain@greekjoin}
\makeatother

\usepackage{geometry}
\geometry{a4paper,margin=25mm,marginparwidth=28mm}
\usepackage{marginnote}
\usepackage{xcolor}
\usepackage[normalem]{ulem}
\setlength{\emergencystretch}{3em}
\usepackage{hyperref}
\hypersetup{colorlinks=true,linkcolor=black,urlcolor=black,pdfborder={0 0 0}}

% --- Batch metadata ---------------------------------------------------------
% Taken as they stand by the rendering script, which shows them at the head of
% the reading view. They fix what the file claims to cover. The macro names are
% the parent project's, so that the scripts stay shared; \folder here means the
% volume's slug — `fr-9115` — and \pages counts Gallica views.

\newcommand{\folder}[1]{\gdef\grothFolder{#1}}
\newcommand{\batch}[1]{\gdef\grothBatch{#1}}
\newcommand{\pages}[2]{\gdef\grothFirst{#1}\gdef\grothLast{#2}}
\newcommand{\foldertitle}[1]{\gdef\grothFoldertitle{#1}}

% The holder's own shelfmark, printed as they write it: « Français 9115 ».
\newcommand{\shelfmark}[1]{\gdef\grothShelfmark{#1}}

% The dating, copied verbatim from the holder's catalogue — never inferred
% here. « XIXe siècle » is the BnF's; a pass that dates a leaf from its content
% says so in a \note{}, as its own inference.
\newcommand{\dating}[1]{\gdef\grothDating{#1}}

% The status notice, printed in the title block and repeated as a diagonal
% watermark on every page of the PDF. The manuscripts are in the public
% domain; what this line declares is not a right but a quality — an
% unverified first pass by a machine. The skills write the line; a file
% without it gets no watermark, so leaving it out is visible in the output.
\newcommand{\watermark}[1]{\gdef\grothWatermark{#1}}

\gdef\grothFolder{?}\gdef\grothBatch{}\gdef\grothFirst{?}\gdef\grothLast{?}%
\gdef\grothFoldertitle{}\gdef\grothShelfmark{}\gdef\grothDating{}\gdef\grothWatermark{}

% --- The résumé -------------------------------------------------------------
\newenvironment{resume}
  {\small\setlength{\parskip}{0.45em}}
  {\par\medskip\hrule\medskip}

% --- Keywords ---------------------------------------------------------------
% In English, closing the modernised reading's résumé: the modern vocabulary
% under which to search for what the leaves construct. The manifest extracts
% them to tag the volume — this line is the single source of the tags.
\newcommand{\keywords}[1]{\par\smallskip\noindent
  {\footnotesize\textsc{Keywords} --- \textit{#1}}\par}

% --- The critical apparatus -------------------------------------------------

% The Gallica view number, in the margin. This is the anchor that lets a
% disagreement over a formula be settled by looking at *one* image rather than
% a volume of 750. The site uses it to turn the facsimile as the transcription
% is scrolled. A view is one scanned image: usually one side of a leaf.
\newcommand{\page}[1]{%
  \marginnote{\footnotesize\color{gray!70}\textsf{v.\,#1}}%
  \label{page:#1}\ignorespaces}

% The BnF's pencilled foliation, where the leaf carries it and a pass can read
% it: « 198r ». This is what the literature cites — « ff. 198r–208v » — and
% the one line that turns a citation into a view. Recorded, never inferred:
% a folio a pass cannot read is not written.
\newcommand{\folio}[1]{%
  \marginnote{\footnotesize\color{gray!50}\textsf{f.\,#1}}\ignorespaces}

% The modernised reading's coarser counterpart to \page{}: which views a
% section is drawn from.
\newcommand{\pagerange}[2]{%
  \marginnote{\footnotesize\color{gray!70}\textsf{v.\,#1--#2}}%
  \ignorespaces}

% Illegible: never guessed.
\newcommand{\ill}{\textcolor{red!60!black}{[\,\ldots\,]}}

% A doubtful reading: the word is given, and flagged as uncertain.
\newcommand{\uncertain}[1]{\uline{#1}}

% An editorial addition — a missing letter, an implied word.
\newcommand{\add}[1]{\textcolor{blue!50!black}{[#1]}}

% Struck out by the writer. What was crossed out is part of the document.
\newcommand{\struck}[1]{\sout{#1}}

% The transcriber's note, on the state of the page or the mathematics.
\newcommand{\note}[1]{\par\smallskip{\small\color{gray!60!black}\textit{[#1]}}\par\smallskip}

% A marginal note of hers, distinct from the body of the text.
\newcommand{\marginal}[1]{\par\smallskip\noindent\hspace{1em}%
  {\small\color{gray!60!black}#1}\par\smallskip}

% --- Title ------------------------------------------------------------------
% The title block is French, like the documents themselves.

\AddToHook{shipout/background}{%
  \ifx\grothWatermark\empty\else
    \begin{tikzpicture}[remember picture, overlay]
      \node[rotate=52, scale=3.4, align=center, text=gray!18,
            font=\sffamily\bfseries]
        at (current page.center) {\grothWatermark};
    \end{tikzpicture}%
  \fi}

\AtBeginDocument{%
  \begin{center}
    {\large\bfseries Manuscrits de mathématiciens (Gallica) ---
      \ifx\grothShelfmark\empty\grothFolder\else\grothShelfmark\fi}\\[2pt]
    {\itshape\grothFoldertitle}\\[4pt]
    {\small vues \grothFirst--\grothLast
      \ifx\grothBatch\empty\else\ (lot \grothBatch)\fi}\\[2pt]
    \ifx\grothDating\empty\else
      {\small\color{gray!60!black}Datation du catalogue : \grothDating}\\[2pt]
    \fi
    \ifx\grothWatermark\empty\else
      {\small\bfseries\color{red!45!gray}\grothWatermark}\\[2pt]
    \fi
    {\footnotesize\color{gray!60!black}Transcription établie d'après les images IIIF de Gallica.
     Source gallica.bnf.fr / Bibliothèque nationale de France.\\
     Les lectures incertaines sont soulignées, les illisibles notées [\,\ldots\,],
     les ajouts entre crochets ; \textsf{v.} numérote les vues de Gallica,
     \textsf{f.} la foliotation au crayon de la BnF.}
  \end{center}
  \vspace{1em}}

\endinput


\folder{nal-2338}
\shelfmark{NAL 2338}
\batch{1}
\pages{1}{20}
\dating{}
\watermark{Lecture automatique\\non vérifiée}
\foldertitle{Papiers divers de Mersenne.}

\begin{document}

\page{2}
\note{Premier feuillet d'une lettre latine, d'une cursive italienne régulière, penchée, à grandes boucles. En haut à droite, à l'encre, d'un trait fin, « 1 » (voir l'en-tête) ; au-dessous, dans une accolade, la cote « Lat. nouv. acq. 2338 ». Au pied, « R.c. 8070 (86) » souligné et le cachet ovale « ACQ. 8070 (LIBRI) » ; au milieu de la page, le cachet rond « BIBLIOTHÈQUE NATIONALE MANUSCRITS » avec « R.F. ».}

M. Mers:\textsuperscript{no} Clariss\textsuperscript{mo}, et Celeberrimo Viro S.P.
\note{« Mers:\textsuperscript{no} » : « Mersenno » abrégé ; « Clariss\textsuperscript{mo} » est écrit avec une grande capitale C bouclée et les lettres finales en exposant.}

Iam non in legendis epistolis tuis Vir Cl\textsuperscript{me} de obscuritate queror, sed de breuitate. adeò enim et eruditionis, et beneuolentiæ erga me tuæ plenæ sunt, ut quantumuis longæ, breuissimæ tamen uideri mihi debeant. Effecimus insuper uterque, tu elegantiori characterum delineatione, ego repetitâ sæpius eorumdem lectionê, ut cursim æquè ac Gallici Typographi, omissâ omni hesitatione, scriptionem tuam perlegam. Ex me quæris an Galileus quidquam reliquerit de ui percussionis: penè nihil; omnesq; illas, et quidem plurimas multumq; uigilatas contemplationes, quas circa hoc subiectum in intellectu habebat, plurimiq; faciebat oculatus ille cæcus, secum tulit. Exceptis tamen duobus experimentis mechanicis, quorum utrumq; inferre uidetur Percussionis uim infinitam esse debere. Duo hæc tantum ex \uncertain{optatissimis} illis cogitatis supersunt; quia Galileus aliquibus amicis suis, inter quos et mihi illa significauerat tamquam primas causas, quæ animum suum impulerant ad considerandum huiusmodi subiectum, atq; in ipso infinitum momentum indagandum. Pulcherrimum sanè problema illud est circa curuam Parabolicam alteri curuæ æqualem: agnosco in hoc titulo subtilissimum ingenium incomparabilis Mathematici Roberuallij, cuius opera utinam uiuente me in lucem prodeant. De hoc autem Problemate, ac de illo numerico quod attulit D. du Verdus nihil habeo quod ad te scribam, quandoquidem circa parabolicum nihil, circa numericum uerò parum admodum immoratus
\note{« optatissimis » est écrit sous le cachet rond de la bibliothèque, qui en couvre les premières lettres. « omnesq; », « multumq; », « plurimiq; », « utrumq; », « atq; » : le q suivi du signe de -que, conservé. Au pied de la page, à droite, le mot « ratus », seul : réclame de la page suivante, qui commence par ce mot (« immo-ratus »).}

\page{3}
\note{Verso du feuillet précédent, même main ; la lettre continue. Aucun numéro sur la page.}

ratus sum, cum non uiderim (ut uera fatear) unde mihi initium speculationis, aditusq; panderetur. De Cycloide, uel Trochoide, habemus etiam solida ante paucos menses inuenta, et Geometris Italis communicata. Attamen si acutiss\textsuperscript{mus} Roberuallius adeò altè naturam huius figuræ penetrauit, quemadmodum ipse refers Vir Cl\textsuperscript{me}, concedo spontè omnium harum inuentionum gloriam Geometræ præstantiss\textsuperscript{o}, et uerè mirabili. Principiò quò ad figuras planas ita definimus. Esto circulus AB, qui Cycloidem Primariam AD describat uti uulgatum est: producaturq; radius CAE in infinitum. Lineas quæ à quolibet puncto radij CA in infinitum producti describuntur Trochoides uel Cycloides secundarias appellabimus. Circulumque uniuscuiusq; proprium genitorem eum uocabimus cuius diameter æqualis sit maximæ suæ Trochoidis altitudini. Iam.
\note{Une tache d'encre couvre le haut de « Cycloide » et la parenthèse qui ferme « (ut uera fatear) ». À droite de « Principiò quò ad figuras \ldots{} uulgatum est », une figure : un cercle de centre C, touchant en A une droite horizontale ; B au sommet du cercle ; le rayon CA prolongé au-dessous de la droite jusqu'à E ; une courbe pointillée part du cercle et monte vers D, à droite.}

Spatium sub qualibet Cycloide linea, et recta eius basi comprehensum ad triangulum eiusdem basis, et altitudinis, est ut peripheria circuli proprij genitoris una cum duplo basis cycloidis, ad duplum basis eiusdem Cycloidis. Vel fortassè melius in rectis lineis, quod idem est, ut patet manifestè, sic

Quodlibet spatium Cycloidale prædictum ad prædictum triangulum, est ut altitudo siue axis eiusdem spatij una cum duplo axis primariæ Cycloidis ad duplum axis eiusdem Primariæ. Ratio, quam quodcumque spatium cycloidale habet ad circulum proprium genitorem, \uncertain{habetur}
\note{« Spatium » et « Quodlibet » commencent en saillie dans la marge de gauche, comme des énoncés. En marge de gauche, à la hauteur de « Cycloidis ad duplum axis », une ligne de points. « prædictum » est écrit avec le p barré de l'abréviation de præ. Au pied, à droite, le mot « in », seul : réclame de la page suivante.}

\page{4}
\note{Même main ; la lettre continue. En haut à droite, à l'encre, d'un trait fin, « 2 » (voir l'en-tête). Les premières lettres des lignes qui commencent en saillie (« Tangens », « In », « Solidum ») passent dans la marge, contre le bord du feuillet.}

in ijsdem præfatis terminis Per conuersionem rationis.
\note{« Per conuersionem rationis » est souligné.}

Tangens ad imperatum punctum cuiuscunq; Cycloidis ita datur. Transeat per datum punctum, peripheria circuli proprij genitoris basim tangentis, quam peripheriam in eodem dato puncto contingat recta linea conueniens cum basi Cycloidis; tum fiat ut altitudo propositæ Cycloidis ad altitudinem Primariæ ita tangens prædicta inter datum punctum, et basim intercepta, ad aliam quandam lineam aptè sumendam in basi ex termino tangentis: deindè ab extremitate huius assumptæ tangens ad imperatum Cycloidis punctum emittatur. Pro ipsa autem basi, sumi potest quælibet ipsi parallela.
\note{« peripheriam » : la première lettre paraît récrite.}

In Primaria Cycloide Tangens procedit à puncto sublimiori genitoris circuli per datum punctum transeuntis \add{basimq; tangentis}. Corollarium \uncertain{ero}.
\note{« basimq; tangentis » est écrit au-dessus de la ligne, appelé par un signe d'insertion après « transeuntis ». Le dernier mot se lit « ero » ; la lecture ne donne pas de sens sûr (« est » ?).}

Solidum quod fit à spatio Cycloidali circa tangentem axi æquidistantem reuolutum, ad cylindrum eiusdem altitudinis, eiusdemq; diametri est subsesquitertium. Huius demonstrationem non habeo, sed inuentum demonstratumq; fuit hoc ab Antonio Nardio patritio Aretino olim Galilei intimo. reliqua mea sunt.

Solidum quod fit à spatio Cyc: circa tangentem basi parallelam reuolutum, est ad cylindrum eiusdem axis, et diametri, \uncertain{subsesquiseptimum}.
\note{le mot est coupé en fin de ligne (« subsesqui\ldots= ») et sa fin, « timum », est écrite au-dessus.}

Solidum quod fit à spatio Cycloidali circa axem reuolutum ad cylindrum eiusdem axis et basis, est ut 11 ad 18.

Solidum circa axem ad solidum circa basim ineffabilem rationem habet: componitur enim ex ratione 44 ad 45. et ex ratione circuli alicuius ad quadratum circumscriptum.

Solidum circa basim ad cylindrum eiusdem axis et diametri est ut 5 ad \uncertain{8}. dabimusq; demonstrationem petentibus; \uncertain{audiuimus}, et fatemur demonstratum fuisse hoc à Cl\textsuperscript{mo} Viro Roberuallio ante nos.
\note{le second chiffre de « 5 ad 8 » est bouclé en haut et peut se lire 6 ; il ressemble au 8 de « 11 ad 18 » deux lignes plus haut. Au pied, à droite, « Centrum », seul : réclame de la page suivante.}

\page{5}
\note{Verso du feuillet « 2 », même main ; la lettre continue. En haut à gauche, une petite marque pâle, qui paraît être le « 2 » du recto vu par transparence ; aucun numéro sur la page.}

Centrum grauitatis integræ Cycloidis axem ita diuidit ut pars ad uerticem terminata sit ad reliquam ut 7 ad 5.

Centrum uerò semicycloidis est in linea axi æquidistante, quæ integræ totius Cycloidis basim ita secat ut partes sint 19. et 8. siue, quæ basim semicycloidis ita secat ut pars ad curuam terminata sit ad reliquam ut 16 ad 11.

Clariss: Donius epistolam à te propriam accepit, et meam perlegit. itaque spero ad ea quæ continentur in utraq; responsurum. Doctissimum æquè, atq; amabilissimum Dominum D. du Verdus uix salutaui. Ipse enim mihi reddidit epistolam tuam, qua iure de silentio meo querebaris Vir Clarissime; tum uno, aut altero tantum, breuiq; colloquio, non nisi pauca simul conferre potuimus circa communia studia Mathematica. Miror quod ille ne uerbum quidem fecerit de misterio illarum tangentium, quas ipse significas. Videbo num per litteras aliquid possim ab ipso impetrare. De libello meo duo opuscula impressa sunt. reliqua duo lentius adhuc procedunt. occupatissimus enim sum tum in officina Vitraria ad usum Telescopij, tum in ædificandis in dies machinis ad tentanda uaria experimenta physica, utrumque iussu Ser\textsuperscript{mi} Domini Magn: Ducis. Nouum opus tuum Cogitata Physicomathematica auidissimè expectamus: ego autem præcipuè, qui hucusque multò plura didici ex lectione unius tui libri Vir Clar\textsuperscript{me}, quam
\note{« de » (après « iure ») est récrit d'un trait lourd sur un premier tracé. « Donius », « du Verdus » : les noms sont écrits ainsi. Au pied, à droite, « plurium », seul : réclame de la page suivante.}

\page{6}
\note{Même main ; fin de la lettre, qui n'occupe que le tiers supérieur de la page, le reste blanc. En haut à droite, à l'encre, d'un trait fin, « 3 » (voir l'en-tête).}

plurium aliorum. adeò eruditione semper atq; omni doctrinarum genere abundant. Æternum certè debeo Clarissimo Donio nostro, quòd primus omnium ipse me dignum putauit, qui uiderem opera tua apud ipsum existentia. Vale diu in publicum bonum Vir eximiè; meque ama, atq; hanc nimiam, et incomptam loquacitatem excusa.

Dat: Flor: Kal: Maij: anno 1644

Addictiss\textsuperscript{mus} Seruus

Euang\textsuperscript{a} Torricellius.
\note{Date et signature autographes en apparence de la même main que la lettre. Le 1 de « 1644 » est une boucle ouverte, lue 1 par le sens. « Euang » est suivi d'une petite boucle en exposant, lue « a » (Euangelista) ; « Torricellius » est souligné d'un trait qui se prolonge. « Addictiss\textsuperscript{mus} » : « mus » est écrit au-dessus de la ligne. La lettre ne porte pas d'adresse sur cette page ; son destinataire n'est nommé qu'en tête, « M. Mers:\textsuperscript{no} » (vue 2).}

\page{7}
\note{Verso du dernier feuillet de la lettre, qui sert d'adresse ; le texte du recto (vue 6) y transparaît à l'envers. Plis de pliage visibles. L'adresse, en haut à droite de la page :}

Celeberrimo Viro P. Marino Mersenno

Lut. Paris.
\note{entre les deux lignes, un grand paraphe bouclé, vertical, non lu. « Lut. Paris. » : Lutetiam Parisiorum, abrégé. Deux petites marques pâles aux coins supérieurs, sans valeur de numéro.}

\page{8}
\note{Une seconde lettre, d'une autre main que celle des vues 2--6 : une italique droite, posée et calligraphiée, presque sans rature, qui garnit les fins de ligne de longs traits de remplissage (non transcrits). En haut à droite, à l'encre, d'un trait fin, « 4 » (voir l'en-tête). Au milieu de la page, le cachet rond de la bibliothèque avec « R.F. ».}

Doctissimo et celeberrimo Viro P. M. Mersenno Euang\textsuperscript{ta}: Torricellius S.
\note{l'intitulé nomme l'auteur ; la main de cette lettre n'est pas celle de la lettre signée « Euang\textsuperscript{a} Torricellius » des vues 2--6. Si elle est une copie ou la main d'un secrétaire, la page ne le dit pas.}

Seris epistolis damus serum responsum: aduenerunt enim litteræ P. V. Clarissimiq; Viri Roberuallij, quamquam Kal. Ian. datæ sint, prope finem Mensis Martij; eæq; mihi redditæ sunt ꝑ Ill\textsuperscript{mos} Viros Comitem Ferdinandum Bardi, et Andream Arrighetti. De mora vero tam seri mei responsi vix ineunte Iulio redditi, non est quod me excusem; inertia enim mea notior est quam vt possit occultari. Excepi satis frequentes querelas circa vitrum illud quod ex meis P. V. voluit habere. Ad hæc primùm respondeo; si vitrum illud idem sit quod ego Romam misi, neque a quopiam (quod non credo) in tot itineribus commutatum sit, optimum esse; neq; in tota Gallia, vniuersisq; regionibus transalpinis vllum æque bonum reperiri censeo, dumodo e manibus meis, vel saltem Francisci Fontanæ emissum non sit. Non habemus quidem oculos Enoc, \& Eliæ, similesque somniantium ineptias, at habemus certissimam demonstrationem nusquam gentium fieri posse vitra meis perfectiora. æque ꝑfecta fieri possibile est, facientq; Pauci, quos æquus amabit Iupiter, et qui Conicorum, refractionumq; doctrinam æque possidebunt, at non meliora. Notum iam est hoc ꝑ vniuersam fere Europam vnde mea vitra expetuntur. At P. V. ne quæso ex quocunq; vitro meo argumentum sumat de reliquis omnibus meis: nam singula eiusdem notæ non euadunt. Vix enim decima quæque ꝑfectissima fiunt ob materiæ imperfectionem, quæ difficillime, nec nisi peracto opere dignosci potest; reliqua partim penitus mala sunt, partim mediocria. Nescio an ornatissimus iuuenis M. Angelus Ricci orauerit P. V. vt vitrum illud ad nos remitteret, pecuniamq; suam haberet: certe ego enixe eum rogaueram vt hoc faceret; neque pecuniam vllam vnquam accepissem (quam ex Florentia
\note{L'initiale S de « Seris » est une grande capitale ornée, en forme de §. « ꝑ » : le p barré de l'abréviation de per. « P. V. » : Paternitas Vestra. La phrase se poursuit à la page suivante.}

\page{9}
\note{Verso du feuillet « 4 », même main posée ; la lettre continue. Aucun numéro sur la page. Au-delà du bord droit du feuillet, sur le papier plus large du feuillet suivant, « Pagina 317 » et « Pagina 327 versu 10 » : ces renvois sont ceux de la vue 10, où ils sont donnés.}

Romam ad eum retro miseram) nisi ipse certum me fecisset P. Vestram tandem nescio quo pacto animaduertisse vitrum non esse penitus contemnendum. Mira scribit P. V. de vitro Gassendiano, quod ipse habuit à Galileo. Vitra nostra, \& præcipuè illud quod P. V. habet, sine dubio meliora sunt quam Vitrum illud electissimum, quod Galileus in delicijs habebat, \& præ oculis suis amabat. Nescio an existimandum sit Galileum ex duobus vitris, melius misisse ad Cl. Gassendum, peius verò sibi ipsi reseruauisse. sed illud absurdum omnino est Tubum Gassendianum quarta parte breuiorem ijsdem adhibitis concauis, maius obiectum representare. Hoc enim idem prorsus est ac si P. V dixisset, Positis duobus triangulis isoscelibus, \& æquiangulis, quod minora habet crura maiorem basim habuisse. Scripsit etiam P. V reperisse quoddam Vitrum Hippoliti Francini amici mei quod nostro præstantius erat. Quodnam fuerit illud vitrum triginta ducatis, hoc est vndecim pistolis emptum optime scio. At iam satis superq; notum est vitra Fontanæ præstantiora Francinianis reperiri, mea verò reliqu omnibus meliora esse. Longitudo tubi argumentum est verum arcanum de figura vitrorum me reperisse. Præteritam hiemem impendi iussu Ser\textsuperscript{mi} M. Ducis in ꝑficiendis Telescopijs, vsq; ad longitudinem decem passuum Geometricorum. Talia verò in Cælum non possunt conuerti sine maximis molibus; tubum fecimus ex malo triremis, scisso prius bifariam, deinde excauato, et iterum compacto; sed tam grauis est, tam longus, tam incomodus, \& fere numquam rectus, vt penè nullius vsus reperiatur. Sequenti hieme quando labor resumi poterit, non ero monendus me debere concauum P. V
\note{« Gassendum » : ainsi écrit à la première occurrence, sans i. « reliqu omnibus » : un blanc suit « reliqu », la fin du mot (reliquis ?) n'est pas écrite. « Longitudinem » commence en saillie dans la marge, après un trait. La phrase se poursuit à la page suivante.}

\page{10}
\note{Même main posée ; la lettre continue. En haut à droite, à l'encre, d'un trait fin, « 5 » (voir l'en-tête). La page est écrite aux trois quarts ; le bas est blanc.}

quod quidem à me factum sit, et ad vitrum 4 brachiorum iustam habeat debitamq; proportionem. Ill\textsuperscript{mus} Vir Io: Baptista Balianus misit ad me nomine P. V quoddam examen Quadraturæ circuli à Batano quodam Mathæmatico editæ. Examen vidi atq; illud vt videretur misi Romam ad amicos meos, breue quidem est, at multò breuius esse poterat, si pro demonstratione quadam numerica, quam præmittit, adhibuisset Propositionem 3 libri VI Euclidis. sed de hoc examine nihil amplius recordor. Habeat hic P. V quadraturam ꝑfectissimam à Mathematico Romano Societatis Iesu demonstratam nuperrime in lib.\textsuperscript{o} cuius titulus Ars Magna lucis, et Vmbræ. Quadratura talis est
\note{« Mathæmatico » : ainsi écrit, avec æ. « Ars Magna lucis, et Vmbræ » est écrit en plus grands caractères, au début d'une ligne en retrait.}

\marginal{Pagina 317}
Si radius AB, arcusque quadrantis AC vterq; bifariam secetur in D, et E, producta DE secabit ex tangente rectam CF æqualem arcui CE.
\note{l'énoncé est écrit en retrait, les lettres de la figure en capitales droites ; aucune figure n'est tracée. Le renvoi « Pagina 317 », en marge de gauche, est de la même main.}

\marginal{Pagina 327 versu 10}
Habet præterea idem Geometra aliquid quod et vestris Geometris, quibus nihil nouum est, fortasse nouum erit. hoc est Quadratum circulo isoperimetrum eidem circulo æquale esse.

Ad Clar\textsuperscript{mu}: Virum Roberuallium mittimus spirales quasdam (vt credo) nostras, quarum linea æqualis ostenditur lineæ rectæ.

Ad Ill\textsuperscript{mu}: D. De Fermat: Hyperbolas (fortasse meas) mitterem, sed vereor ne ineptias meas penitus dedignetur videre, qui iam diu assueuit nobilissimo infinitarum parabolarum Theoremati. Certè quod ad numerum attinet pauciores non erunt hyperbolæ nræ.
\note{« nræ » est surmonté du trait d'abréviation (nostræ). Le bas de la page est blanc ; la lettre se termine au verso (vue 11).}

\page{11}
\note{Verso du feuillet « 5 », même main posée ; fin de la lettre commencée à la vue 8. Le haut de la page seulement est écrit ; le texte du recto transparaît, à l'envers, dans le reste de la page. Aucun numéro sur la page.}

Tibi verò V. Clar\textsuperscript{me}: Carollariolum mitto ex ipsis hyperbolis deductum. Quadratura quædam est, quarum centenas, immò infinitas poteram mittere, nisi vidissem satis superque esse vnam, vt statim omnes emergant. Iamque Vale, meque obsequentissimum seruum tuum ama.
\note{« Carollariolum » : ainsi écrit, avec a ; lire Corollariolum. Aucun corollaire n'est écrit sur la page : ce qui était joint à la lettre n'y est pas.}

\uncertain{D.} Flor: die 7 Iulij 1646. Euangelista Torricellius.
\note{la date est précédée d'une grande capitale bouclée, lue D (Datum ?) avec doute. La signature est de la même main que le corps de la lettre et de même module : elle n'a pas l'aspect de la signature de la vue 6.}

Quocunq; tempore P. V iusserit pecuniam suam habebit, vitrumq; ad nos remittet si ita volet
\note{post-scriptum, séparé de la signature par un blanc, sans ponctuation finale. Le sujet de « habebit » et de « remittet » n'est pas nommé ; la lettre parle plus haut de « M. Angelus Ricci » (vue 8).}

\page{12}
\note{Une troisième lettre, de la main rapide des vues 2--6. En haut à droite, à l'encre, d'un trait fin, « 6 » (voir l'en-tête). Au milieu de la page, le cachet rond de la bibliothèque avec « R.F. ».}

Cl\textsuperscript{mo}: Viro P. Marino Mersenno

Euang\textsuperscript{ta} Torricellius S. D.
\note{« Mersenno » est d'une encre plus pâle que le reste de la ligne ; qu'il ait été ajouté après coup ou d'une autre main, la page ne permet pas de le dire.}

Pluribus epistolis tuis unicum responsum dabo Vir Clar\textsuperscript{me} doctissimeq;: Doctissimum uerò appello quinto quodam modo, cuius rei rationem aperiam, quotiescunque docueris ipse, quarè me incomparabilem Geometram, sed quarto quodam modo, limitata silicet appellatione dixeris in una ex tuis epistolis. dicis enim rationem apodicticam te daturum, si tamen ratio ulla demonstratiua dari potest illius tui dicti, quod ex omni parte falsum est. Sed omissis iocis ueniamus ad seria. Sæpè iam excusauisti, Vir omnifariam, et quocunq; modo uolueris doctissime, tarditatem responsi quod à Cl\textsuperscript{mo} Viro Roberuallio uestro debetur litteris meis iam ab ineunte mense Augusto ad ipsum datis. Sine quæso ipsum quantumcunq; libuerit cunctari. Neque enim ipsi respondere potero usque ad futuram æstatem, si tamen unquam potero. iam enim litteræ nostræ in immensum crescunt, et epistolæ degenerant in uolumina. At mihi quælibet iactura temporis grauis est, et multa habeo potiora, quæ maturanda essent, et tamen omnia prætermitto.
\note{« epistolis », à la première ligne, est à demi effacé. « silicet » : ainsi écrit. « prætermitto », comme plus bas « præferam » et « prætulissem », est écrit avec le p lié à un e cédillé.}

De centro illo percussionis, de quo iam non semel scripsisti, nihil sollicitus sum. Problema enim non arridet animo meo. Tubum quem dedi Cl: Viro Bullialdo æquè perfectum existimarem, sed octaua ferè parte breuiorem isto tuo. Verum cum ipse centies iam scripseris istum tuum nihil penitus ualere, superest ergo ut ego tubum Bullialdi præferam, quem nunquam prætulissem, nisi me certiorem fecisses de imperfectione istius, quod tam sæpe reperisti facto experimento omnibus aliorum uitris inferius. Veniet cum hac epistola uitrum cauum quod iam diu promiseram. Expectaui
\note{Dans cette main, le double s ressemble à un n (« fecisses », « prætulissem »). Au pied, à droite, « diu mul: », seul : réclame de la page suivante.}

\page{13}
\note{Verso du feuillet « 6 », même main rapide ; la lettre continue. Aucun numéro sur la page. À droite de l'ouverture paraît le bord du feuillet suivant (vue 14), dont la page est donnée à sa vue.}

diu multumq; reditum alicuius ex uestris in Galliam, ut ad te illud perferret, sed frustra: nondum enim contigit mihi alloqui aliquem ad uos redeuntem. Sed quid tandem proderit uitrum cauum mittere, si præcipuum illud conuexum tam malum est? Vitra mea iam undique toto terrarum orbe expetuntur, quamquam adeo imperfecta tibi uisa sint V. Cl\textsuperscript{me}. Hoc unum moneo, uitrum illud quod ipse habes hoc requirere, silicet ut eius pars detecta, uel nihil penitus differat ab hoc circulo, uel si differet paulò admodum maior sit ad obiecta obscura, minor uerò ad clariora. Cæterum quicunq; iudicat tubum tuum ideò meliorem Gassendiniano teloscopio quia maiora refert obiecta cum longior sit, pace illius dictum uelim, fallitur. Melior enim est qui purissimam uisionem offert, et distinctiora exhibet obiecta præsertim ea quæ ualdè minuta sunt. Legi iam sæpe inauguralem epistolam Clar\textsuperscript{mi} Gassendi magna cum animi uoluptate, maiore cum admiratione. Summas ago tibi gratias Vir doctiss\textsuperscript{me} qui me dignum tam pretioso munere indicauisti. Quotiescunq; ad me litteras misisti pro eruditiss\textsuperscript{mo} Viro Io: Baptista Donio, ipsas a me statim accepit, et summoperè caras habuisse testatus est. Tubis meis longissimis nihil adhuc noui deteximus in Cælo, præter fascias Iouiales, quæ ipsum Iouis globum, tanquam terrestres nostræ zonæ, præcingunt. Ipse enim ferè nunquam in Cælum aspicio ob inopportunitatem ædium
\note{À droite de « ab hoc circulo », un cercle tracé à la plume, d'environ quatre lignes de diamètre : le « circulus » du texte, la mesure de l'ouverture que la lettre prescrit. « maiora » : le a final est écrit au-dessus d'une lettre noircie. « silicet », « Gassendiniano », « teloscopio » : ainsi écrits. Au pied, à droite, « quas », seul : réclame de la page suivante.}

\page{14}
\note{Même main rapide ; la lettre continue. En haut à droite, à l'encre, d'un trait fin, « 7 » (voir l'en-tête). À gauche de l'ouverture, le bord du verso précédent (vue 13), avec son cercle.}

quas inhabito. Spero iam Roberuallium uestrum non amplius dicturum impossibiles Spirales meas quæ infinitas reuolutiones peragunt antequam ad centrum perueniant: credo enim ipsum iam reperisse demonstrationes omnes, cum illarum quas Geometricas, et penitus meas appello, tum etiam illarum quas ex Archimedis definitione in infinitas Algebræ dignitates propagata, deriuo. Idem etiam spero de Problemate infinitarum hyperbolarum. Quid enim non credam de tali ingenio, cui tamen (ignoscas mihi hoc Vir Clar\textsuperscript{me}) ipse detrahis aliquid, dum tantum intempestiuis laudibus extollere contendis. Scripsisti iam per uniuersam ferè Italiam litteris Genuam, Bononiam, Romam, et Florentiam missis, mira atq; inaudita de epistola illa quam adornat Roberuallius. Vocas enim et mirabilem, et mirabilium inuentorum plenam, et inimitabilem, et maximam in Geometria confidentiam præseferentem, atq; his similia. Quisnam tantæ expectationi ab omni parte respondere possit nisi unus Roberuallius, et iudicio tuo! Siquis (ut ipse scripsisti Vir Clar\textsuperscript{me}) ex Typographis urbis uoluerit libellos meos iterum typis committere, gratum faciet si curabit ut plurimi errores qui in p\textsuperscript{ma} editione irrepserunt, tollantur. ipse quod prædictis libellis addam in præsens nihil polliceor. Habeo quidem plura Geometrica quorum editionem aliquando fortasse expediam si per ætatem licebit, sed paratis prius his, ut cecinit ille, quorum indiget usus. Proximè præteritis mensibus Autumni incidi in \uncertain{Problema}
\note{« antequam » est écrit « anteq\ldots » avec la finale en boucle. Dans le dernier mot, le b n'est pas distinct (« Prolema »). Au pied, à droite, « quoddam », seul : réclame de la page suivante.}

\page{15}
\note{Verso du feuillet « 7 », même main rapide ; fin de la lettre commencée à la vue 12. Aucun numéro sur la page ; le bas de la page est blanc, le recto y transparaît.}

quoddam propositum, ut ego audiui, ab Ill\textsuperscript{mo} Viro de Fermat. Iubebat enim datis tribus punctis aliud reperire ex quo tres eductæ ad data tria puncta sint minima quantitas. Construxi Problema, demonstraui, determinauiq;, nam propositum fuerat sine determinatione solutio non una est: alia enim per doctrinam solidorum procedit; alia atq; alia sine locis solidis per pura plana rem omnem absoluit. Si uolueris demonstrationes habebis uel ex me, uel ex Caualerio, uel Magiotto, uel Renerio nostro; cum uarijs enim amicis illas contul\add{i} quamquam facilis admodum contemplatio uideatur. Doceas quæso, num huiusmodi Problematis solutio apud uos in uulgus exierit, an apud Autorem hactenus lateat. Audio ab ingeniosiss\textsuperscript{mo} amico nro M. Ang.\textsuperscript{o} Riccio utriq; nostrûm destinatum fuisse librum nouissimè recusum Niceronis uestri, ab ipso Autore dum inter uiuos ageret. deuinces me singulari beneficio si curaueris supremam hanc uoluntatem optimi, celeberrimiq; uiri adimpleri. Interim certò scias Vir Clar\textsuperscript{me} te nunquam erraturum donec inter magis certas et immutabiles possessiones tuas adnumerabis Torricellium. Vale.
\note{« nam propositum fuerat sine determinatione solutio non una est » : ainsi écrit, sans ponctuation entre « determinatione » et « solutio ». « contul » est au bord droit du feuillet ; la dernière lettre manque. « nro » est surmonté du trait d'abréviation (nostro). La lettre ne porte pas d'autre signature que le nom écrit dans la dernière phrase.}

Florentiæ Kalendis Februarij anno 1647.

\page{16}
\note{Une quatrième pièce, d'une autre main encore : une cursive penchée à gros traits, aux finales en boucle, qui écrit souvent e pour un e final d'infinitif surmonté d'un signe. En haut à droite, à l'encre, d'un trait fin, « 8 » (voir l'en-tête). Au milieu de la page, le cachet rond de la bibliothèque avec « R.F. ». Rien sur la page ne nomme celui qui écrit « ego scripsi » ; la lettre qu'il dit avoir reçue est citée au style direct sous son intitulé.}

Doctissime, et Celeberrime P. Mersenne

Oro P. Vestram ut secum ipsa recordari uelit quando ego scripsi Centrum grauitatis Cycloidis secare axem in ratione 7 ad 5, et solidum circa axem esse ad cylindrum ut 11 ad \uncertain{18}, ipsam mihi in Epistola maximè longa responsum hoc dedisse sub die \uncertain{28} Iunij 1644
\note{le second terme de « 11 ad 18 » se lit aussi 10 : le 8 de cette main est à demi formé. Le jour, « 28 », est écrit d'un 2 bouclé et d'un 8 en forme de \&.}

Incomparabili Geometræ D. Torricellio S.P.D.

Vix credere possis V. Clar. Charissimeq; quantis tuæ nouissimè litteræ accessionibus tuam apud me famam, et æstimationem promouerint. Quid enim illi putem ἀδύνατον qui uel nostrum Geometram Roberuallum inuentione centri grauitatis Cycloidis et illius circa axem solidi * uerbum deest. Reliqua enim inuenit, et demonstrauit. Sed quæ (licet ille parata dicat habere quæ Trochoidis suæ plana spectant ad edendum parata) tamen non debeas infodere tuas circa idem negotium demonstrationes: iuuat enim idem pluribus modis demonstratum inspicere. Ille uerò non solum per indiuisibilia sed more alterius demonstrationis, quam ad te misi, omnia prædicta demonstrauit. qui cum tuas postremas legisset prædictum solidum, et centrum grauitatis tibi debere fatetur, qui primus inuenisti.
\note{« ἀδύνατον » est écrit en lettres grecques. « uerbum deest » est encadré d'un trait, et appelé par un astérisque au-dessus de « solidi » : une remarque de celui qui copie, qui signale un mot manquant dans ce qu'il copie. En marge de gauche, encadré de même, « Parata bis legitur », qui relève la répétition de « parata » dans la parenthèse. « qui cum tuas postremas legisset \ldots{} qui primus inuenisti » est souligné d'un trait continu. Le cachet de la bibliothèque couvre en partie « non debeas » et « iuuat enim ».}
\marginal{Parata bis legitur}

Rogamus tamen an centrum grauitatis solidorum Trochoidicor\textsuperscript{m} habeas, quæ numerasti, ut habes centrum grauitatis plani Cycloidalis; \add{et} cur nam dicas te habere demonstrationem solidi circa basim ut 5 ad \uncertain{8}. Numquid et aliorum habes?
\note{« et » est écrit au-dessus de la ligne, avant « cur ». Le 8 de « 5 ad 8 » a la même forme ouverte que plus haut.}

Omittimus maximam, et longissimam epistolæ partem in qua plures, et etiam clariores huiusmodi confessiones leguntur. Deinde
\note{« Omittimus \ldots{} » : celui qui écrit déclare n'avoir pas recopié la plus grande partie de la lettre qu'il cite. « Deinde » termine la page et ouvre la suite.}

\page{17}
\note{Verso du feuillet « 8 », même main à gros traits ; la pièce continue. Aucun numéro sur la page. Le bord droit du feuillet passe sous le feuillet suivant : les dernières lettres de plusieurs lignes y sont cachées, restituées ci-dessous, comme ajouts de l'éditeur, seulement où le sens est sûr.}

uersus finem iterum hæc habet P. V\textsuperscript{ra}.

Dubitat noster Roberuallus an Mechanicè tantum centra grauitatis Cycloidis, et semicycloidis inueneris, quæ Geometricè falsa suspicatur. docebis num istius rei demonstrationem habeas.
\note{« Dubitat » et « quæ Geometricè falsa suspicatur » sont soulignés.}

Potest ne aliquid clarius desiderari? Postquam ego uidi Cl. Roberuallum suspicari, et P. Vestram a me demonstrationem petere, uix lecta epistola statim misi demonstrationem Centri grauitatis Cycloidis, Solidique circa basim, et quod summoperè doleo in ipsa demonstratione, quæ satis longa erat, misi etia\add{m} demonstrationem meam, et uerè meam, pro methodo, quæ inseru\add{it} ad inueniendum Centrum grauitatis ex dato solido, siue solidum ex dato centro etc. Clariss: tandem Roberuallus in ultima Epistola inquit non solum centrum grauitatis Cycloidis iam diu habuisse, sed etiam methodum meam inuersa tantum Propositione inter sua numerat, quod egerrimè fero. Si enim Centrum grauitatis antequam demonstr\add{a}tionem meam uideret non habebat, quemadmodum certè, imm\uncertain{o} certissimè scio non habuisse (ut P. V., uel ipsemet, uel tandem uniuersa Europa testis esse poterit) sine dubio neque Methodu\add{m} habebat. Nolui primum epistolæ caput ad Clarissimum Roberuallum mittere. satis enim duxi si illud consideraret P. V\textsuperscript{ra}, nam spero ipsam huic meæ iniuriæ obuiam ituram ad ut mihi mea tribuantur; in ipsa enim maximè confido: et
\note{Dans cette main, le t final est une boucle en forme de 8 (« erat », « habebat », « et »). « antequam » est écrit « antiquam » en apparence, le e lié au n. « inseru » et « ad », en fin de ligne, sont au bord caché du feuillet ; « ad » est suivi d'un point ou d'un trait d'union, et la fin du mot (adeo ?) n'est pas visible. La phrase se poursuit à la page suivante.}

\page{18}
\note{Même main à gros traits ; la pièce continue. En haut à droite, à l'encre, d'un trait fin, « 9 » (voir l'en-tête). Le bas de la page est blanc.}

ipsa me protegere debet, quæ à me demonstrationem illam petiuit, et accepit, et quæ semper fuit interpres huius commercij ex parte tantum uestra tam \uncertain{docti} atq; eruditi. Scio etiam \uncertain{eam esse} Cl. Roberuall: humanitatem atq; fidem, eamq; habere ipsum inuentorum suorum copiam, ut statim atq; monitus fuerit à P. Vestra de ratione temporum, de Epistolis datis, et de hoc quod fortasse exciderat ei tot occupationibus distracto, ipsum credam in meam sententiam uenturum. Ne uerò duo prædicta epistolæ capita a me conficta existimentur, neue quidquam siue additum, siue detractum commutatumq; credatur, primùm reminiscentia P. V. fidem facere poterit, deinde auctoritas Clar. Virorum characteres P. V. optimè cognoscentium. Si Ill\textsuperscript{mus} Donius aderat, ab ipso petijssem testimonium de fide mea in describendis capitibus epistolæ prædictæ. Inuoco fidem atq; beneuolentiam P. V. quam maximam censeo, namque uim patior. Si Centrum illud grauitatis sciebat una cum methodo uniuersali Cl. Roberuallus quando ipsi P. V\textsuperscript{ra} meam solam enunciationem ostendit, certè non dixisset hoc ipsum mihi debere, neque me primum Inuentorem P. V\textsuperscript{ra} nominauisset, neque illud falsum Geometricè potuisset suspicari Roberuallius. neque P. V publicè etiam typis edidisset illa Problemata esse mea.
\note{« docti » : la finale peut se lire e (docte) ; « eam esse » : les lettres donnent « eadesse », lues ainsi avec doute. « Roberuall: » est abrégé ; la même page écrit aussi « Roberuallus » et « Roberuallius ».}
\marginal{\uncertain{aduit} tandem licet breui tempore}
\note{note encadrée dans la marge de gauche, à la hauteur de « cognoscentium. Si Ill\textsuperscript{mus} Donius aderat », d'un module plus petit, comme « Parata bis legitur » à la vue 16 ; elle glose « Donius aderat ». Le premier mot, coupé « ad-uit », peut être « adiit » ou « aduenit ».}

Quæso P. V. ignoscas mihi, si fortasse in his multo longius
\note{La phrase se poursuit à la page suivante.}

\page{19}
\note{Verso du feuillet « 9 », même main à gros traits ; fin de la pièce commencée à la vue 16. Aucun numéro sur la page. Le bord droit du feuillet passe sous le feuillet suivant et cache la fin de plusieurs lignes. Le bas de la page est blanc.}

quam oportebat prouectus sum. fateor enim omnia inuenta me\add{a} pro nihilo me habere, et meras nugas cognoscere, at nimiò dolore afficerer et grauiorem contumeliam paterer, si quidqua\add{m} mihi tamquam penitus fatuo, et semimortuo tam manifestè præripi uiderem sineremq;. Incredibile est quanto desiderio expectemus responsum P. V. circa hoc negotium.

Lecta iterum epistola Cl. Roberuallij et obsignata iam mea ad ipsum data, animaduerto me nihil respondisse de solido Cycloidis circa axem. Sed neq; responsum quodpiam dare necesse existimo. Tunc enim quispiam \struck{\ill{}} \add{iure} arguere poterit me quando in paralogismos meos incidet. Habemus apud Archimedem Prop. 2 de Circuli dimensione Circulum ad quad\textsuperscript{m} diametri esse ut 11 ad 14 Quæro ab ipso undenam putet me habuisse rationem quam ad numeros 11 et \uncertain{18} reducebam? si uer\add{ò} ea dicit ut ego demonstrationes iterum \uncertain{ultrò} mittam, fallitur. Prælectionem, siue orationem adeò laudatam doctissimi et eloquentissimi P. Gassendi libentissimè uiderem. Nihil enim ex eius operibus uidi præter uitam D. De Peiresik. Libentissimè etiam scirem ubinam degat hoc temporis Cl. et celeberrimus Des Cartes, nam aliquod commercium cum tan\add{to} Viro ualdè desidero. Atque hic me tibi iterum commendo

Torricellium
\note{« iure » est écrit au-dessus d'un mot noirci et barré. Le 8 de « 11 et 18 » a la forme ouverte relevée à la vue 16 ; il peut se lire 10. « ultrò » : lu « ulerò », avec doute. « putet », « dicit », « degat » : le t final est la boucle en 8 de cette main. « Peiresik » : ainsi écrit. « Torricellium », seul au bas à droite, à la place d'une signature, est le complément de « commendo » ; la pièce n'a ni date ni autre souscription.}

\page{20}
\note{Une cinquième pièce, de la main posée des vues 8--11 (italique droite, traits de remplissage en fin de ligne, non transcrits). En haut à droite, à l'encre, d'un trait fin, « 10 » (voir l'en-tête). En haut à gauche, d'une autre main, plus tardive, souligné : « 1 Sept. 1643 » ; cette date n'est pas écrite dans la lettre et n'est pas de la main du copiste. Au milieu de la page, le cachet rond de la bibliothèque avec « R.F. ».}

Clar\textsuperscript{mo} Viro Roberuall\textsuperscript{o}

Euang\textsuperscript{ta} Torricellius S.P.

\add{A}pertè tecum sine alio interprete, Vir Clar\textsuperscript{me} (quis \add{d}issimulare possit) et quamquam literæ tuæ ad Clar\textsuperscript{m} \ill{}ennum missæ sint, cohibere tamen non possum animi \add{i}mpetum, quin ad te currat, tibique totum se dedicet, tanquam Apollini Geometrarum Fortunatas certè iam existimare debeo \add{nu}gas meas, atque illas non iam amplius nihil facere, quandoquidem digna habita sunt, quæ iudicium tuum subirent, et animaduersionibus tuis nobilitarentur.
\note{La grande initiale de la lettre n'a été qu'esquissée : on n'en voit, dans la marge, que le trait supérieur, en forme de F. Les quatre premières lignes sont en retrait pour la loger, et les lettres qui commencent ces lignes et les deux suivantes, dans la même zone, ont pâli ou manquent : « \add{A}pertè », « \add{d}issimulare », « \add{i}mpetum » (le point du i est visible), « \add{nu}gas » sont restitués par le sens. Au début de la troisième ligne, le mot qui suit « ad Clar\textsuperscript{m} » n'est lisible qu'à partir de « ennum » ; le nom qu'il porte n'est pas restitué.}

Principiò ex me quæris, an centrum grauitatis parabolæ à priori vt inuentum à me proponatur, aut quæratur vt ignotum: erubescerem certè ignotum Theorema inter alias Propositiunculas meas à me demonstratas collocare. ostendimus illud vnica, breuique demonstratione. sed ea occasione admiratus sum fœcunditatem ingenij tui, circa tot parabolas, atque earum solida, non solum Geometricè, sed etiam Mechanicè considerata, et ad mensuram, scientiamque redacta. De hijs nihil ego habeo quod proferam, et fortasse non habebo; siquidem, difficillimæ, nisi fallor, contemplationis censeo huiusmodi Theor; \textsuperscript{te} præterea immorari \uncertain{ul} soleo circa figuras non vulgatas, et circa solida, quæ si noua sint, saltem ab antiquis, et receptis figuris planis ortum non habeant atque hoc ea præcipuè ratione, vt laborum fructus, quando res ex animi voto succedat, communem litteratorum applausum \uncertain{sortiatur}; neque sit qui inuideat figuris a me ipso fabricatis.
\note{« Theor; » finit la ligne ; « te » est écrit au-dessus, après le point-virgule : fin abrégée du mot (Theoremata ?) ou premier mot de la ligne suivante, la page ne le décide pas. « ul », devant « soleo », est un u suivi d'une boucle haute, non résolu. « sortiatur » se lit « cortiatur ». Le cachet de la bibliothèque couvre « etiam » et « redacta » sans les rendre illisibles.}

Mensura Cycloidis, (hoc enim nomine Clar: Galileus appellauit 45 iam ab hinc annis figuram, quæ fortasse tibi nunc Trochois est) mihi se, se vltrò obtulit nl speranti, penè dixi nl quærenti. Illam deinde quinquies diuersis semper principijs
\note{« se, se » : ainsi écrit, pour sese. « nl » : un n suivi d'une haste barrée, l'abréviation de nil (ou de non). La phrase se poursuit au-delà de cette page, à la vue 21.}

\end{document}
